\documentclass[12pt]{article}

% Idioma y codificación
\usepackage[spanish, es-tabla]{babel}       %es-tabla para que se titule "Tabla"
\usepackage[utf8]{inputenc}

% Márgenes
\usepackage[a4paper,top=3cm,bottom=2.5cm,left=3cm,right=3cm]{geometry}

% Comentarios de bloque
\usepackage{verbatim}

% Paquetes de links
\usepackage[hidelinks]{hyperref}    % Permite enlaces
\usepackage{url}                    % redirecciona a la web

% Más opciones para enumeraciones
\usepackage{enumitem}


%------------------------------------------------------------------------

%Paquetes de figuras
\usepackage{caption}
\usepackage{subcaption} % Figuras al lado de otras
\usepackage{float}      % Poner figuras en el sitio indicado H.


% Paquetes de imágenes
\usepackage{graphicx}       % Paquete para añadir imágenes
\usepackage{transparent}    % Para manejar la opacidad de las figuras

% Paquete para usar colores
\usepackage[dvipsnames]{xcolor}
\usepackage{pagecolor}      % Para cambiar el color de la página

% Habilita tamaños de fuente mayores
\usepackage{fix-cm}


%------------------------------------------------------------------------

% Paquetes de matemáticas
\usepackage{amsmath, amsfonts, amssymb, mathrsfs}
\usepackage{cancel}     % Simplificar tachando
\usepackage{polynom}    % Divisiones y Ruffini

\usepackage{pgfplots}   %Representar funciones
\pgfplotsset{compat=1.18}  % Versión 1.18

%Definición de teoremas, etc.
\usepackage{amsthm}
%\swapnumbers   % Intercambia la posición del texto y de la numeración

\theoremstyle{plain}
\newtheorem{teo}{Teorema}[chapter]     % Se resetea en cada chapter
\newtheorem{coro}{Corolario}[teo]           % Se resetea en cada teorema
\newtheorem{prop}[teo]{Proposición}         % Usa el mismo contador que teorema
\newtheorem{lema}[teo]{Lema}                % Usa el mismo contador que teorema

\theoremstyle{remark}
\newtheorem*{observacion}{Observación}

\theoremstyle{definition}
\newtheorem{definicion}{Definición} [chapter]     % Se resetea en cada chapter
\newtheorem*{notacion}{Notacion}
\newtheorem*{ejemplo}{Ejemplo}
\newtheorem*{ejercicio*}{Ejercicio}             % No numerado
\newtheorem{ejercicio}{Ejercicio} [section]     % Se resetea en cada section


% Modificar el formato de la numeración del teorema "ejercicio"
\renewcommand{\theejercicio}{%
  \ifnum\value{section}=0 % Si no se ha iniciado ninguna sección
    \arabic{ejercicio}% Solo mostrar el número de ejercicio
  \else
    \thesection.\arabic{ejercicio}% Mostrar número de sección y número de ejercicio
  \fi
}


% \renewcommand\qedsymbol{$\blacksquare$}         % Cambiar símbolo QED
%------------------------------------------------------------------------

% Paquetes para encabezados
\usepackage{fancyhdr}
\pagestyle{fancy}
\fancyhf{}

\newcommand{\helv}{ % Modificación tamaño de letra
\fontfamily{}\fontsize{12}{12}\selectfont}
\setlength{\headheight}{15pt} % Amplía el tamaño del índice


%\usepackage{lastpage}   % Referenciar última pag   \pageref{LastPage}
\fancyfoot[C]{\thepage}

%------------------------------------------------------------------------

\begin{comment} % Descomentar en book. Comentar en article

% Conseguir que no ponga "Capítulo 1". Sino solo "1."
\makeatletter
\renewcommand{\@makechapterhead}[1]{%
  \vspace*{50\p@}%
  {\parindent \z@ \raggedright \normalfont
    \ifnum \c@secnumdepth >\m@ne
      \huge\bfseries \thechapter.\hspace{1em}\ignorespaces
    \fi
    \interlinepenalty\@M
    \Huge \bfseries #1\par\nobreak
    \vskip 40\p@
  }}
\makeatother
\end{comment}

%------------------------------------------------------------------------
% Personalizar la portada
\usepackage{titling}


%------------------------------------------------------------------------

% Comandos definidos
\newcommand{\bb}[1]{\mathbb{#1}}
\newcommand{\cc}[1]{\mathcal{#1}}
\newcommand{\Hop}{\stackrel{L'H\hat{o}pital}{=}}

\newcommand{\red}[1]{{\color{red}#1}} % Para integrales, destacar los cambios.

% Método de integración
\newcommand{\MetInt}[2]{
    \left[\begin{array}{c}
        #1 \\ #2
    \end{array}\right]
}

%------------------------------------------------------------------------
\newcommand{\T}[0]{\cc{T}}

\begin{document}

    % 1. Foto de fondo
    % 2. Título
    % 3. Encabezado Izquierdo
    % 4. Color de fondo
    % 5. Coord x del titulo
    % 6. Coord y del titulo
    % 7. Fecha

    
    % 1. Foto de fondo
% 2. Título
% 3. Encabezado Izquierdo
% 4. Color de fondo
% 5. Coord x del titulo
% 6. Coord y del titulo
% 7. Fecha

\newcommand{\portada}[7]{

    \portadaBase{#1}{#2}{#3}{#4}{#5}{#6}{#7}
    \portadaBook{#1}{#2}{#3}{#4}{#5}{#6}{#7}
}

\newcommand{\portadaExamen}[7]{

    \portadaBase{#1}{#2}{#3}{#4}{#5}{#6}{#7}
    \portadaArticle{#1}{#2}{#3}{#4}{#5}{#6}{#7}
}




\newcommand{\portadaBase}[7]{

    % Tiene la portada principal y la licencia Creative Commons
    
    % 1. Foto de fondo
    % 2. Título
    % 3. Encabezado Izquierdo
    % 4. Color de fondo
    % 5. Coord x del titulo
    % 6. Coord y del titulo
    % 7. Fecha
    
    
    \thispagestyle{empty}               % Sin encabezado ni pie de página
    \newgeometry{margin=0cm}        % Márgenes nulos para la primera página
    
    
    % Encabezado
    \fancyhead[L]{\helv #3}
    \fancyhead[R]{\helv \nouppercase{\leftmark}}
    
    
    \pagecolor{#4}        % Color de fondo para la portada
    
    \begin{figure}[p]
        \centering
        \transparent{0.3}           % Opacidad del 30% para la imagen
        
        \includegraphics[width=\paperwidth, keepaspectratio]{assets/#1}
    
        \begin{tikzpicture}[remember picture, overlay]
            \node[anchor=north west, text=white, opacity=1, font=\fontsize{60}{90}\selectfont\bfseries\sffamily, align=left] at (#5, #6) {#2};
            
            \node[anchor=south east, text=white, opacity=1, font=\fontsize{12}{18}\selectfont\sffamily, align=right] at (9.7, 3) {\textbf{\href{https://losdeldgiim.github.io/}{Los Del DGIIM}}};
            
            \node[anchor=south east, text=white, opacity=1, font=\fontsize{12}{15}\selectfont\sffamily, align=right] at (9.7, 1.8) {Doble Grado en Ingeniería Informática y Matemáticas\\Universidad de Granada};
        \end{tikzpicture}
    \end{figure}
    
    
    \restoregeometry        % Restaurar márgenes normales para las páginas subsiguientes
    \pagecolor{white}       % Restaurar el color de página
    
    
    \newpage
    \thispagestyle{empty}               % Sin encabezado ni pie de página
    \begin{tikzpicture}[remember picture, overlay]
        \node[anchor=south west, inner sep=3cm] at (current page.south west) {
            \begin{minipage}{0.5\paperwidth}
                \href{https://creativecommons.org/licenses/by-nc-nd/4.0/}{
                    \includegraphics[height=2cm]{assets/Licencia.png}
                }\vspace{1cm}\\
                Esta obra está bajo una
                \href{https://creativecommons.org/licenses/by-nc-nd/4.0/}{
                    Licencia Creative Commons Atribución-NoComercial-SinDerivadas 4.0 Internacional (CC BY-NC-ND 4.0).
                }\\
    
                Eres libre de compartir y redistribuir el contenido de esta obra en cualquier medio o formato, siempre y cuando des el crédito adecuado a los autores originales y no persigas fines comerciales. 
            \end{minipage}
        };
    \end{tikzpicture}
    
    
    
    % 1. Foto de fondo
    % 2. Título
    % 3. Encabezado Izquierdo
    % 4. Color de fondo
    % 5. Coord x del titulo
    % 6. Coord y del titulo
    % 7. Fecha


}


\newcommand{\portadaBook}[7]{

    % 1. Foto de fondo
    % 2. Título
    % 3. Encabezado Izquierdo
    % 4. Color de fondo
    % 5. Coord x del titulo
    % 6. Coord y del titulo
    % 7. Fecha

    % Personaliza el formato del título
    \pretitle{\begin{center}\bfseries\fontsize{42}{56}\selectfont}
    \posttitle{\par\end{center}\vspace{2em}}
    
    % Personaliza el formato del autor
    \preauthor{\begin{center}\Large}
    \postauthor{\par\end{center}\vfill}
    
    % Personaliza el formato de la fecha
    \predate{\begin{center}\huge}
    \postdate{\par\end{center}\vspace{2em}}
    
    \title{#2}
    \author{\href{https://losdeldgiim.github.io/}{Los Del DGIIM}}
    \date{Granada, #7}
    \maketitle
    
    \tableofcontents
}




\newcommand{\portadaArticle}[7]{

    % 1. Foto de fondo
    % 2. Título
    % 3. Encabezado Izquierdo
    % 4. Color de fondo
    % 5. Coord x del titulo
    % 6. Coord y del titulo
    % 7. Fecha

    % Personaliza el formato del título
    \pretitle{\begin{center}\bfseries\fontsize{42}{56}\selectfont}
    \posttitle{\par\end{center}\vspace{2em}}
    
    % Personaliza el formato del autor
    \preauthor{\begin{center}\Large}
    \postauthor{\par\end{center}\vspace{3em}}
    
    % Personaliza el formato de la fecha
    \predate{\begin{center}\huge}
    \postdate{\par\end{center}\vspace{5em}}
    
    \title{#2}
    \author{\href{https://losdeldgiim.github.io/}{Los Del DGIIM}}
    \date{Granada, #7}
    \thispagestyle{empty}               % Sin encabezado ni pie de página
    \maketitle
    \vfill
}
    \portadaExamen{ffccA4.jpg}{Topología I\\Examen VIII}{Topología I. Examen VIII}{MidnightBlue}{-8}{28}{2023-2024}{Arturo Olivares Martos}

    \begin{description}
        \item[Asignatura] Topología I.
        \item[Curso Académico] 2023-24.
        \item[Grado] doble Grado en Ingeniería Informática y Matemáticas.
        \item[Grupo] A.
        \item[Profesor] Jose Antonio Gálvez López\footnote{El examen lo pone el departamento}.
        \item[Descripción] Convocatoria Ordinaria.
        \item[Fecha] 18 de enero de 2023.
        \item[Duración] 3 horas.
    
    \end{description}
    \newpage
    
    \begin{ejercicio}[5 puntos]
        Sea $\| \cdot \|$ la norma euclídea usual en $\mathbb{R}^2$. Se consideran los conjuntos
        $X = \{(x,y) \in \mathbb{R}^2 \mid \| (x,y) \| \geq 2\}$ y $S = \{(x,y) \in \mathbb{R}^2 \mid \| (x,y) \| = 2\}$.
        Para cada $p = (x,y) \in X$, se considera la siguiente familia:
        \begin{itemize}
          \item Si $p \in S$, $\beta_p = \{B((0,0), r) \cup \{p\} \mid 0 < r < 1\}$.
          \item Si $p \notin S$, $\beta_p = \{B(p,r) \mid 0 < r < 2 - \| p \|\}$.
        \end{itemize}
        Responde razonadamente a las siguientes preguntas:
        \begin{enumerate}
          \item(1.5 puntos) Demuestra que existe una única topología $\T$ en $X$ tal que las familias anteriores forman una base de entornos de cada punto. En los siguientes apartados, se considera el espacio topológico $(X,\T)$.
          \item(1 punto) Dado el subconjunto $A = \{(x,y) \in X : x > 0\}$, calcula el interior y la frontera de $A$.
          \item(0.5 puntos) Estudia si $(X,\T)$ Hausdorff o no.
          \item(1 punto) Estudia si $(X,\T)$ es conexo o no.
          \item(1 punto) Estudia si $(X,\T)$ es compacto o no.
        \end{enumerate}

    \end{ejercicio}

    \begin{ejercicio}[2.5 puntos]
        Elige una de las siguientes preguntas (2a o 2b):
        \begin{enumerate}
          \item[2a.] Da una definición de subespacio compacto en un espacio topológico y prueba las siguientes afirmaciones:
          \begin{enumerate}
            \item Un subespacio cerrado de un espacio topológico compacto es compacto.
            \item Todo subespacio compacto de un espacio topológico Hausdorff es cerrado.
          \end{enumerate}

          \item[2b.] Razona si son verdaderas o falsas las siguientes afirmaciones:
          \begin{enumerate}
            \item Todo entorno de un punto de un espacio topológico es abierto.
            \item Si $(X, \T)$ es un espacio topológico entonces la aplicación identidad dada por $Id:~(X, \T) \to (X, \T_{\text{CF}})$ es continua si y solo si $(X,\T)$ es T1. Aquí, $\T_{\text{CF}}$ denota la topología cofinita.
            \item Si $f : (X, \T) \to (Y, \T')$ es una aplicación biyectiva entre espacios compactos y T2, entonces $f$ es continua si, y solo si, $f^{-1}$ es continua.
          \end{enumerate}
        \end{enumerate}
    \end{ejercicio}
        
    \begin{ejercicio}[2.5 puntos]
        En $X = \mathbb{R} \times [-1,1]$ consideramos la relación de equivalencia dada por:
        \[(x_1,y_1)\cc{R}(x_2,y_2) \iff \left\{
            \begin{array}{ll}
            (x_1, y_1) = (x_2, y_2), & \text{o bien} \\
            x_1, x_2 > 1, & \text{o bien} \\
            x_1, x_2 < -1.
            \end{array}
        \right.\]
        Si sobre $X$ consideramos la topología usual inducida de $\mathbb{R}^2$, estudiese si:
        \begin{enumerate}
            \item La proyección canónica $p : X \to X/\cc{R}$ es una aplicación abierta,
            \item $X/\cc{R}$ es T2,
            \item $X/\cc{R}$ es compacto,
            \item $X/\cc{R}$ es conexo.
        \end{enumerate}
    \end{ejercicio}
\end{document}